\documentclass[11pt]{article}
\usepackage[margin=1in]{geometry}
\usepackage{amsmath,amssymb,amsthm,hyperref}
\hypersetup{hidelinks}
\newtheorem{proposition}{Proposition}
\title{Disproof of Conjecture 00000003711\\Resistance sums and spanning trees of a triangle}
\author{gaochengzhecpu}\date{}
\begin{document}\maketitle
\noindent\textbf{Claim addressed.} The conjecture identifies the sum of the resistance-distance quadratic forms with the number of spanning trees. The connected simple graph $K_3$, with unit conductance on each edge, disproves this identification under both unordered-pair and ordered-pair conventions. The individual pseudoinverse formula for resistance is not disputed.

\begin{proposition}
For the unit triangle, every pair of distinct vertices has effective resistance $2/3$. The sum over unordered vertex pairs is $2$, the sum over ordered pairs is $4$, and the number of spanning trees is $3$.
\end{proposition}
\begin{proof}
Let $I$ be the $3\times3$ identity and $J$ the all-ones matrix. The graph Laplacian and its Moore--Penrose inverse are
\[
 L=3I-J=\begin{pmatrix}2&-1&-1\\-1&2&-1\\-1&-1&2\end{pmatrix},
 \qquad Q=\frac13I-\frac19J.
\]
Indeed $J^2=3J$, so $LQ=QL=I-J/3$. Direct multiplication gives $LQL=L$, $QLQ=Q$, and both $LQ$ and $QL$ are symmetric. These are the four defining Moore--Penrose equations.

For vertices $i,j$, write $b=e_i-e_j$. Then $Jb=0$, and hence
\[
 v=Qb=\frac13b,\qquad Lv=b.
\]
This is the voltage solving Kirchhoff's law for a unit current entering at $i$ and leaving at $j$. Its terminal drop equals the quadratic form:
\[
 R_{ij}=v_i-v_j=b^TQb=
 \begin{cases}0&i=j,\\2/3&i\ne j.\end{cases}
\]
The drop does not depend on the voltage representative. If $Lw=0$, subtracting the first two equations gives $w_0=w_1$, and either equation then gives $w_2=w_0$. Thus two solutions differ by a constant.

There are three unordered pairs and six ordered distinct pairs. Their sums are consequently $3(2/3)=2$ and $6(2/3)=4$. A spanning tree of $K_3$ consists of exactly two of its three edges. Each such choice is connected and acyclic, and there are exactly three choices. Therefore neither resistance sum equals the spanning-tree count $3$.
\end{proof}

\noindent\textbf{Scope.} This is a connected, unweighted, loopless simple graph; no missing connectedness hypothesis is exploited. Including diagonal pairs adds zero. The source's assertion about the sum fails with either usual pair-counting convention. No assertion is made against the usual matrix-tree theorem or the individual effective-resistance formula.

\section*{Lean correspondence and reproduction}
The project defines $K_3$ as Mathlib's complete \texttt{SimpleGraph (Fin 3)}. Its actual \texttt{lapMatrix} is proved equal to $L$. The four Moore--Penrose equations for the real matrix $Q$ are proved entry by entry. The current, voltage and quadratic form are genuine matrix-vector and dot products. In addition, \texttt{electrical\_resistance} proves that every solution of the actual Kirchhoff equation has the claimed terminal drop.

Spanning trees are subgraphs on the full vertex set satisfying Mathlib's actual \texttt{IsTree}. They are placed in bijection with the three-edge masks having exactly two edges. The proof uses the connectedness and edge-count characterization of finite trees, with finite connectedness checked by the library's walk-based decision procedure. Thus the number $3$ is a derived cardinality, not a count assigned by definition.

The final theorems \texttt{unordered\_conjecture\_false} and \texttt{ordered\_conjecture\_false} compare the actual quadratic-form sums with that cardinality.

Inside \texttt{lean/}, run \texttt{lake build}, followed by
\begin{center}\texttt{lake env lean -DwarningAsError=true Main.lean}.\end{center}
Lean is pinned to 4.19.0 and Mathlib to
\begin{center}\small\texttt{c44e0c8ee63ca166450922a373c7409c5d26b00b}.\end{center}
The dependency manifest pins every transitive dependency. The supplementary \texttt{verify.py} independently checks the matrix equations and tree enumeration with exact rational arithmetic. There are no proof gaps, custom axioms or native computation shortcuts. Actual fresh-build logs, axiom dependencies and artifact hashes are in \texttt{verification/}. Author and delegated-agent review records are separate; no external independent review is claimed.

\begin{thebibliography}{9}
\bibitem{source} The Last Math Competition,
\href{https://github.com/The-Last-Math-Competition/The-Last-Math-Competition/blob/main/conjectures/00000003711.md}{Conjecture 00000003711 (original bilingual source)}.
\end{thebibliography}
\end{document}
